
\documentclass{article}
\usepackage{colortbl}
\usepackage{makecell}
\usepackage{multirow}
\usepackage{supertabular}

\begin{document}

\newcounter{utterance}

\centering \large Interaction Transcript for game `eqbench', experiment `eq\_bench', episode 27 with reflexion{-}llama{-}3.1{-}8b{-}instruct.
\vspace{24pt}

{ \footnotesize  \setcounter{utterance}{1}
\setlength{\tabcolsep}{0pt}
\begin{supertabular}{c@{$\;$}|p{.15\linewidth}@{}p{.15\linewidth}p{.15\linewidth}p{.15\linewidth}p{.15\linewidth}p{.15\linewidth}}
    \# & \multicolumn{2}{c}{Player} && \multicolumn{2}{c}{Game Master} \\
    \hline

    \theutterance \stepcounter{utterance}  
    & & & \multicolumn{4}{p{0.6\linewidth}}{
        \cellcolor[rgb]{0.9,0.9,0.9}{
            \makecell[{{p{\linewidth}}}]{
                \texttt{\tiny{[P1$\langle$GM]}}
                \texttt{You have attempted this same task before and did not succeed. These are the plans you wrote after those attempts.} \\
\texttt{Trial \#1: New plan: I got stuck in a situation where I didn't accurately assess Sophia's emotional response. Upon re{-}reading the dialogue, I realize that Sophia's tone is more disapproving and critical than Jamal's defensive and assertive tone. I should have considered the power dynamics and the context of the conversation. Next time, I will pay closer attention to the language and tone used by each character, and consider the potential emotional responses based on their perspectives and relationships. I will also consider the fact that Sophia is trying to enforce project guidelines, which might make her feel more threatened or annoyed than excited or apathetic. I will adjust my scores accordingly, with Annoyed and Threatened being more likely emotions for Sophia.} \\
\\ 
\texttt{Trial \#2: New plan: I got stuck in a situation where I didn't accurately assess Sophia's emotional response. Upon re{-}reading the dialogue, I realize that Sophia's tone is more disapproving and critical than Jamal's defensive and assertive tone. I should have considered the power dynamics and the context of the conversation. Next time, I will pay closer attention to the language and tone used by each character, and consider the potential emotional responses based on their perspectives and relationships. I will also consider the fact that Sophia is trying to enforce project guidelines, which might make her feel more threatened or annoyed than excited or apathetic. I will adjust my scores accordingly, with Annoyed and Threatened being more likely emotions for Sophia. Specifically, I will:} \\
\\ 
\texttt{{-} Increase the score for Annoyed to 9, as Sophia's tone is disapproving and critical, and she is trying to enforce project guidelines that Jamal is not following.} \\
\texttt{{-} Increase the score for Threatened to 8, as Sophia is trying to maintain control and authority in the conversation, and Jamal's unconventional design drafts might be perceived as a challenge to her authority.} \\
\texttt{{-} Decrease the score for Excited to 0, as there is no indication in the dialogue that Sophia is feeling excited or enthusiastic about the design drafts.} \\
\texttt{{-} Decrease the score for Apathetic to 0, as Sophia's tone and language suggest that she is actively engaged in the conversation and has strong opinions about the design drafts.} \\
\\ 
\texttt{Annoyed: 9} \\
\texttt{Excited: 0} \\
\texttt{Threatened: 8} \\
\texttt{Apathetic: 0} \\
\\ 
\texttt{Trial \#3: New plan: I got stuck in a situation where I didn't accurately assess Sophia's emotional response. Upon re{-}reading the dialogue, I realize that Sophia's tone is more disapproving and critical than Jamal's defensive and assertive tone. I should have considered the power dynamics and the context of the conversation. Next time, I will pay closer attention to the language and tone used by each character, and consider the potential emotional responses based on their perspectives and relationships. I will also consider the fact that Sophia is trying to enforce project guidelines, which might make her feel more threatened or annoyed than excited or apathetic. I will adjust my scores accordingly, with Annoyed and Threatened being more likely emotions for Sophia.} \\
\\ 
\texttt{To improve my assessment, I will:} \\
\\ 
\texttt{{-} Increase the score for Annoyed to 9, as Sophia's tone is disapproving and critical, and she is trying to enforce project guidelines that Jamal is not following.} \\
\texttt{{-} Increase the score for Threatened to 9, as Sophia is trying to maintain control and authority in the conversation, and Jamal's unconventional design drafts might be perceived as a challenge to her authority.} \\
\texttt{{-} Decrease the score for Excited to 0, as there is no indication in the dialogue that Sophia is feeling excited or enthusiastic about the design drafts.} \\
\texttt{{-} Decrease the score for Apathetic to 0, as Sophia's tone and language suggest that she is actively engaged in the conversation and has strong opinions about the design drafts.} \\
\\ 
\texttt{Annoyed: 9} \\
\texttt{Excited: 0} \\
\texttt{Threatened: 9} \\
\texttt{Apathetic: 0} \\
\\ 
\\ 
\\ 
\texttt{Your task is to predict the likely emotional responses of a character in this dialogue:} \\
\\ 
\texttt{Sophia: Jamal, I've been looking at your design drafts. They're... unconventional.} \\
\texttt{Jamal: That's one way to put it, Soph. I prefer to call it 'innovative'.} \\
\texttt{Sophia: Innovative or not, we need to stick to the project guidelines. We can't afford to take risks.} \\
\texttt{Jamal: Risks? That's what creativity is all about, Sophia. We need to break the mold.} \\
\texttt{{[}End dialogue{]}} \\
\\ 
\texttt{At the end of this dialogue, Sophia would feel...} \\
\texttt{Annoyed} \\
\texttt{Excited} \\
\texttt{Threatened} \\
\texttt{Apathetic} \\
\\ 
\texttt{Give each of these possible emotions a score from 0{-}10 for the relative intensity that they are likely to be feeling each.} \\
\\ 
\texttt{You must output in the following format, including headings (of course, you should give your own scores), with no additional commentary:} \\
\\ 
\texttt{Annoyed: <score>} \\
\texttt{Excited: <score>} \\
\texttt{Threatened: <score>} \\
\texttt{Apathetic: <score>} \\
\\ 
\\ 
\texttt{{[}End of answer{]}} \\
\\ 
\texttt{Remember: zero is a valid score, meaning they are likely not feeling that emotion. You must score at least one emotion > 0.} \\
\\ 
\texttt{Your answer:} \\
            }
        }
    }
     \\ \\

    \theutterance \stepcounter{utterance}  
    & \multicolumn{4}{p{0.6\linewidth}}{
        \cellcolor[rgb]{0.561,0.737,0.561}{
            \makecell[{{p{\linewidth}}}]{
                \texttt{\tiny{[P1$\rangle$GM]}}
                \texttt{Annoyed: 8} \\
\texttt{Excited: 0} \\
\texttt{Threatened: 6} \\
\texttt{Apathetic: 0} \\
            }
        }
    }
    & & \\ \\

    \theutterance \stepcounter{utterance}  
    & & & \multicolumn{2}{p{0.3\linewidth}}{
        \cellcolor[rgb]{0.9,0.9,0.9}{
            \makecell[{{p{\linewidth}}}]{
                \texttt{\tiny{[GM$|$GM]}}
                \texttt{Annoyed: 8} \\
\texttt{Excited: 0} \\
\texttt{Threatened: 6} \\
\texttt{Apathetic: 0} \\
            }
        }
    }
    & & \\ \\

    \theutterance \stepcounter{utterance}  
    & & & \multicolumn{2}{p{0.3\linewidth}}{
        \cellcolor[rgb]{0.9,0.9,0.9}{
            \makecell[{{p{\linewidth}}}]{
                \texttt{\tiny{[GM$|$GM]}}
                \texttt{Annoyed: 7} \\
\texttt{Excited: 0} \\
\texttt{Threatened: 4} \\
\texttt{Apathetic: 0} \\
            }
        }
    }
    & & \\ \\

    \theutterance \stepcounter{utterance}  
    & & & \multicolumn{2}{p{0.3\linewidth}}{
        \cellcolor[rgb]{0.9,0.9,0.9}{
            \makecell[{{p{\linewidth}}}]{
                \texttt{\tiny{[GM$|$GM]}}
                \texttt{game\_result = LOSE} \\
            }
        }
    }
    & & \\ \\

    \theutterance \stepcounter{utterance}  
    & & \multicolumn{4}{p{0.6\linewidth}}{
        \cellcolor[rgb]{0.95,0.95,0.95}{
            \makecell[{{p{\linewidth}}}]{
                \texttt{\tiny{[P1$\leftrightarrow$P]}}
                \texttt{Post{-}episode Reflexion (not sent to GM):} \\
\texttt{New plan: I got stuck in a situation where I didn't accurately assess Sophia's emotional response. Upon re{-}reading the dialogue, I realize that Sophia's tone is more disapproving and critical than Jamal's defensive and assertive tone. I should have considered the power dynamics and the context of the conversation. Next time, I will pay closer attention to the language and tone used by each character, and consider the potential emotional responses based on their perspectives and relationships. I will also consider the fact that Sophia is trying to enforce project guidelines, which might make her feel more threatened or annoyed than excited or apathetic. I will adjust my scores accordingly, with Annoyed and Threatened being more likely emotions for Sophia.} \\
\\ 
\texttt{To improve my assessment, I will:} \\
\\ 
\texttt{{-} Increase the score for Annoyed to 10, as Sophia's tone is disapproving and critical, and she is trying to enforce project guidelines that Jamal is not following.} \\
\texttt{{-} Increase the score for Threatened to 10, as Sophia is trying to maintain control and authority in the conversation, and Jamal's unconventional design drafts might be perceived as a challenge to her authority.} \\
\texttt{{-} Decrease the score for Excited to 0, as there is no indication in the dialogue that Sophia is feeling excited or enthusiastic about the design drafts.} \\
\texttt{{-} Decrease the score for Apathetic to 0, as Sophia's tone and language suggest that she is actively engaged in the conversation and has strong opinions about the design drafts.} \\
\\ 
\texttt{Annoyed: 10} \\
\texttt{Excited: 0} \\
\texttt{Threatened: 10} \\
\texttt{Apathetic: 0} \\
            }
        }
    }
    & \\ \\

\end{supertabular}
}

\end{document}
